% =========================================================================
% Paper KR-FP-B : Mass gap for 4D pure Yang-Mills via Bakry-Emery on the
%                 fundamental modular domain : a conditional reduction.
%
% Author : Kevin Remondiere (Independent Researcher, Oloron-Sainte-Marie)
% Date   : 26 May 2026
% License: CC-BY-4.0
% Target : Letters in Mathematical Physics (LMP) / Annals of Mathematics
% =========================================================================

\documentclass[11pt,a4paper]{amsart}

\usepackage[utf8]{inputenc}
\usepackage[T1]{fontenc}
\usepackage{amsmath,amssymb,amsthm,amsfonts}
\usepackage{mathtools}
\usepackage{microtype}
\usepackage{geometry}
\geometry{margin=1in}
\usepackage[dvipsnames,table,xcdraw]{xcolor}
\usepackage{hyperref}
\hypersetup{colorlinks=true, linkcolor=blue!50!black, citecolor=blue!50!black, urlcolor=blue!50!black}
\usepackage{booktabs}
\usepackage{enumitem}

\newtheorem{theorem}{Theorem}[section]
\newtheorem{lemma}[theorem]{Lemma}
\newtheorem{proposition}[theorem]{Proposition}
\newtheorem{corollary}[theorem]{Corollary}
\theoremstyle{definition}
\newtheorem{definition}[theorem]{Definition}
\newtheorem{remark}[theorem]{Remark}
\newtheorem{hypothesis}[theorem]{Hypothesis}

% -- Macros --------------------------------------------------------------
\newcommand{\R}{\mathbb{R}}
\newcommand{\Z}{\mathbb{Z}}
\newcommand{\gfrak}{\mathfrak{g}}
\newcommand{\hfrak}{\mathfrak{h}}
\newcommand{\su}{\mathfrak{su}}
\newcommand{\norm}[1]{\left\lVert#1\right\rVert}
\newcommand{\abs}[1]{\left\lvert#1\right\rvert}
\newcommand{\inner}[2]{\langle #1, #2\rangle}
\newcommand{\Ric}{\mathrm{Ric}}
\newcommand{\Tr}{\mathrm{Tr}}
\newcommand{\ad}{\mathrm{ad}}
\newcommand{\Coul}{\mathrm{Coul}}
\newcommand{\dvol}{\, d\mathrm{vol}}
\DeclareMathOperator{\spec}{spec}
\DeclareMathOperator{\rk}{rk}
\DeclareMathOperator{\Hess}{Hess}
\DeclareMathOperator{\Ent}{Ent}

% IMPORTANT : everywhere below, \kFP denotes the Faddeev--Popov / Kostant
% invariant 1/(2|\Phi^+(G)|), NOT the entanglement-entropy area-law
% prefactor (which is a distinct object, see RemondiereKRFP3_2026).
\newcommand{\kFP}{\kappa_{\mathrm{FP}}}
\newcommand{\cinf}{c_{\infty}}

% =========================================================================
\title[Mass gap via Bakry--\'Emery on $\Lambda$]
{Mass gap for 4D pure Yang--Mills via Bakry--\'Emery on the
 fundamental modular domain: a conditional reduction}

\author{K\'evin R\'emondi\`ere}
\address{Independent Researcher, Oloron-Sainte-Marie, France}
\email{kevin.remondiere@gmail.com}
\thanks{ORCID:
\href{https://orcid.org/0009-0008-2443-7166}{0009-0008-2443-7166}.
This work is released under a
\href{https://creativecommons.org/licenses/by/4.0/}{Creative Commons
Attribution 4.0 International License (CC-BY-4.0)}.}
\date{26 May 2026}
\subjclass[2020]{Primary 81T13; Secondary 35P15, 47D07, 53C21, 58J50, 60J60.}
\keywords{Yang--Mills mass gap, Faddeev--Popov operator, fundamental modular domain,
Bakry--\'Emery, logarithmic Sobolev inequality, Polchinski equation, cluster
expansion, Otto--Villani.}

\begin{document}

\begin{abstract}
We articulate a conditional reduction of the 4D pure $SU(N)$ Yang--Mills mass
gap problem to two named external inputs:
(i)~the spectral bound KR--FP--3 of the companion work~\cite{RemondiereKRFP3_2026}
on the Faddeev--Popov operator $M[A] = d_A^\dagger d_A$ over the closure
$\overline{\Lambda}_{S_0}$ of an action-bounded slice of the fundamental
modular domain $\Lambda$ of Singer~\cite{Singer1978} and
Dell'Antonio--Zwanziger~\cite{DellAntonioZwanziger1991}, viz.\
$\lambda_{\min}(M[A]) \geq m_0^2 (1-\kFP)$ with
$\kFP = 1/(2|\Phi^+(G)|)$;
(ii)~a logarithmic Sobolev inequality with constant uniform in the lattice
regularisation, in the spirit of the Bauerschmidt--Bodineau--Dagallier
Polchinski-equation program~\cite{BauerschmidtBodineauDagallier2024,
BauerschmidtDagallier2024}, applied to the Wilson measure of 4D pure
$SU(N)$ lattice gauge theory.
Under (i) and (ii), the Babelon--Viallet O'Neill formula~\cite{BabelonViallet1981}
translates the Faddeev--Popov bound into a Bakry--\'Emery curvature--dimension
estimate $\Ric_{\mathcal A/\mathcal G} \geq (1-\kFP) m_0^2 \cdot g$
on $\overline{\Lambda}_{S_0}$, the Bakry--\'Emery theorem~\cite{BakryEmery1985}
yields a uniform LSI on the symbolic Markov semigroup, and the
Rothaus~\cite{Rothaus1981} / Otto--Villani~\cite{OttoVillani2000} chain
delivers a positive spectral gap that survives the continuum descent via
the Fukushima--Oshima--Takeda Dirichlet-form
machinery~\cite{FukushimaOshimaTakeda1994}.
For $G=SU(3)$ and $D=4$, the explicit conditional bound is
$\Delta \geq \tfrac{5}{6}\, m_0^2 / \cinf(4) > 0$.
We emphasise that the reduction is \emph{conditional}: the three Faddeev--Popov
hypotheses (H1)~generic-vanishing, (H2)~compact-manifold Sobolev, (H3)~measurable
Cartan selection of~\cite{RemondiereKRFP3_2026} remain open, and the
infinite-volume descent inherits the open problem of cluster-expansion control
for the 4D $SU(N)$ Wilson measure. The contribution of the present note is
strictly the geometric/analytic reduction; the construction itself remains
out of reach.
\end{abstract}

\maketitle

% =========================================================================
\section{Introduction}\label{sec:intro}

The Yang--Mills mass gap problem, one of the seven Clay Millennium Prize
problems~\cite{JaffeWitten2000}, asks for a mathematically rigorous
construction of the 4D quantum Yang--Mills field theory with compact simple
gauge group $G$ together with a proof that the spectrum of the Hamiltonian
satisfies $\spec(H)\setminus\{0\}\subset[\Delta,\infty)$ for some
$\Delta>0$. Recent decades have seen two principal analytic routes:
the renormalisation-group / cluster-expansion approach pioneered by
Ba\l aban~\cite{Balaban1985} and revisited in modern probabilistic form
by Chandra--Chevyrev--Hairer--Shen in
3D~\cite{ChandraChevyrevHairerShen2024} and by Chatterjee for
$SU(2)$-Higgs~\cite{Chatterjee2024}, and the variational / dynamical
approach via stochastic quantisation and Polchinski-type
equations~\cite{BauerschmidtBodineauDagallier2024,BauerschmidtDagallier2024,
CaoParkSheffield2024,CaoNissimSheffield2025}.

\paragraph{Geometric route.} A third, classically motivated route consists of
working on the orbit space $\mathcal A/\mathcal G$ directly, exploiting its
Riemannian structure (O'Neill formula of
Babelon--Viallet~\cite{BabelonViallet1981}, Mitter--Viallet~\cite{MitterViallet1981})
and the bottom of the spectrum of the Faddeev--Popov operator
$M[A] = d_A^\dagger d_A$ as a curvature lower bound. The principal obstruction
to this geometric route has been the absence of a uniform spectral bound on
$M[A]$ over the relevant domain inside $\mathcal A/\mathcal G$, due to the
Gribov ambiguity~\cite{Gribov1978}: any naive Coulomb-gauge slice intersects
$\partial\Omega$ where $M[A]$ is non-invertible.

The companion work~\cite{RemondiereKRFP3_2026} (henceforth ``KR--FP--3'')
articulates a conditional resolution of this obstruction on the fundamental
modular domain $\Lambda$ of Singer~\cite{Singer1978} via three conditional
named hypotheses (H1)--(H3), a Birman--Schwinger / Aubin--Talenti
bound~\cite{Aubin1976,Talenti1976}, and a Lie-algebraic reduction encoded
in the Kostant identity~\cite{Kostant1959}. The output is the conditional
spectral bound
\begin{equation}\label{eq:KRFP3-bound}
\lambda_{\min}(M[A]) \;\geq\; m_0^2\,(1-\kFP), \qquad
\kFP \,:=\, \tfrac{1}{2|\Phi^+(G)|},
\quad A \in \overline\Lambda_{S_0}.
\end{equation}
For $G = SU(3)$, $\kFP = 1/6$ and the bound reads
$\lambda_{\min} \geq \tfrac{5}{6}\, m_0^2$.

\paragraph{Purpose of this Letter.} We articulate the
\emph{conditional reduction}
$$
\eqref{eq:KRFP3-bound}\;+\;
\text{uniform LSI \`a la BBD~\cite{BauerschmidtBodineauDagallier2024}}
\;\Longrightarrow\;
\text{mass gap }\Delta \geq (1-\kFP) m_0^2 / \cinf(D) > 0,
$$
making explicit each transitional step in the chain
$$
\underbrace{M[A]\geq(1-\kFP)m_0^2}_{\text{KR--FP--3}}
\;\xrightarrow{\text{Babelon--Viallet O'Neill}}\;
\underbrace{\Ric_{\mathcal A/\mathcal G}\geq(1-\kFP) m_0^2\,g}_{\text{KR--FP--A}}
\;\xrightarrow{\text{Bakry--\'Emery}}\;
\text{CD}((1-\kFP) m_0^2,\infty)
$$
$$
\;\xrightarrow{\text{Otto--Villani}}\;
\text{uniform LSI}
\;\xrightarrow{\text{Rothaus}}\;
\text{spectral gap}
\;\xrightarrow{\text{Fukushima--Oshima--Takeda}}\;
\text{continuum gap.}
$$
The present Letter does not claim the mass gap. It is a structural
contribution that organises an existing geometric route into a single
chain of named theorems, identifying explicitly which open problems
(KR--FP--3 hypotheses H1--H3 plus the Bauerschmidt--Bodineau--Dagallier
cluster expansion for the 4D $SU(N)$ Wilson measure) remain.

\paragraph{Main theorem (informal).} Conditional on the validity of
\eqref{eq:KRFP3-bound} on $\overline\Lambda_{S_0}$ \emph{and} on a uniform
logarithmic Sobolev inequality with constant $\cinf(D)$ for the lattice
$SU(N)$ Wilson measure in the spirit of~\cite{BauerschmidtBodineauDagallier2024,
BauerschmidtDagallier2024}, the continuum 4D $SU(N)$ pure Yang--Mills mass gap
satisfies
\begin{equation}\label{eq:main-informal}
\Delta \;\geq\; \frac{(1-\kFP)\, m_0^2}{\cinf(D)} \;>\; 0.
\end{equation}
For $G = SU(3)$, $D = 4$, this reads $\Delta \geq \tfrac{5}{6}\,m_0^2/\cinf(4)$.
Section~\ref{sec:thm} states this rigorously as
Theorem~\ref{thm:main} (Theorem~KR--FP--B).

\paragraph{Honest scope.} The author wishes to underscore two boundaries of
the present contribution:
\begin{enumerate}[label=(\alph*),leftmargin=2em]
\item It is \emph{not} a self-contained proof of the Clay mass gap. The
three hypotheses (H1)--(H3) of \cite{RemondiereKRFP3_2026} remain open
analytic problems. The strongest among them, (H1)~generic-vanishing of
minimising eigenfunctions on the closure of the fundamental modular domain,
is supported by lattice numerics but not yet rigorously proved.
\item The uniform LSI input (ii) is, for 4D $SU(N)$ pure gauge,
\emph{not yet established}. The Bauerschmidt--Bodineau--Dagallier program
delivers it for $\varphi^4_2,\varphi^4_3$ scalar
measures~\cite{BauerschmidtDagallier2024}; its adaptation to non-Abelian
gauge measures is the standard cluster-expansion problem. Our role is
to identify it as the \emph{single} remaining external input once
\cite{RemondiereKRFP3_2026} is fully discharged.
\end{enumerate}
The reduction therefore constitutes a partial bypass of the
Ba\l aban--Bauerschmidt cluster expansion \emph{for the Faddeev--Popov
spectral lower bound}, but still requires that program for the measure
descent.

\paragraph{Structure of the Letter.} Section~\ref{sec:prelim} fixes
notation. Section~\ref{sec:KRFPchain} recalls the KR--FP--1/2/3 chain
of~\cite{RemondiereKRFP3_2026} and states Corollary~KR--FP--A.
Section~\ref{sec:bakry} formalises the Bakry--\'Emery step and the passage
to a logarithmic Sobolev inequality. Section~\ref{sec:descent} sets out
the continuum descent via Fukushima--Oshima--Takeda Dirichlet forms,
conditional on the BBD uniform LSI. Section~\ref{sec:thm} states the
conditional theorem. Section~\ref{sec:honest} lists the open hypotheses
and outlines a 2--4 weeks formalisation roadmap conditional on a Bauerschmidt
collaboration. Section~\ref{sec:conclusion} concludes.

% =========================================================================
\section{Preliminaries}\label{sec:prelim}

\subsection{Yang--Mills setup}
Let $(M, g_M)$ be a compact oriented Riemannian four-manifold without
boundary, $P \to M$ a smooth principal $G$-bundle with $G$ a compact simple
Lie group, $\gfrak = \mathrm{Lie}(G)$ with normalised Killing form
$-\Tr(\ad_X\ad_Y) = \langle X, Y\rangle_\gfrak$. The space of connections
$\mathcal A = \Omega^1(M;\mathrm{ad}\,P)$ carries the natural $L^2$
metric. The gauge group $\mathcal G = \Gamma(\mathrm{Ad}\,P)$ acts on
$\mathcal A$ by $A \mapsto A^g = g^{-1}Ag + g^{-1}dg$.

\subsection{Coulomb gauge and the Faddeev--Popov operator}
In Coulomb gauge $\partial^\mu A_\mu = 0$, the Faddeev--Popov operator
acting on $\mathrm{ad}\,P$-valued functions is
$$
M[A]\eta \;=\; d_A^\dagger d_A \eta \;=\; -D^\mu[A]\,D_\mu[A]\,\eta
\;=\; -\Delta\eta + g\,K(A)\eta + g^2 K_2(A)\eta,
$$
with $K(A)\eta = -2\,[A^\mu,\partial_\mu\eta]$ (Coulomb) and
$K_2(A)\eta = -[A^\mu,[A_\mu,\eta]]$.

\subsection{Gribov region and the fundamental modular domain}
Gribov~\cite{Gribov1978}:
$\Omega = \{A : \partial^\mu A_\mu = 0,\ M[A] > 0\}$.
Singer~\cite{Singer1978}:
$\Lambda = \{A \in \Omega : \norm{A}_{L^2}^2 = \min_{g\in\mathcal G}
\norm{A^g}_{L^2}^2\}$, the \emph{fundamental modular domain}.
Dell'Antonio--Zwanziger~\cite{DellAntonioZwanziger1991}:
$\overline{\Lambda}\subset\overline{\Omega}$, and the inversion of $M[A]$
near $\partial\Omega$ is the central analytic difficulty.

\subsection{Action slice}
For $S_0 < \infty$, set
$\overline{\Lambda}_{S_0} := \{A \in \overline{\Lambda} : S_\mathrm{YM}[A]
\leq S_0\}$ with $S_\mathrm{YM}[A] = \tfrac14 \int_M \abs{F[A]}^2\dvol$.
By Uhlenbeck's compactness theorem~\cite{Uhlenbeck1982} adapted to
$\overline{\Lambda}$, this set is weakly precompact in
$H^1_{\Coul}(M; \Omega^1\otimes\mathrm{ad}\,P)$
(see~\cite[Lemma~5.1 of the companion]{RemondiereKRFP3_2026}).

\subsection{Bakry--\'Emery curvature--dimension condition}
For a symmetric Markov semigroup $(P_t)_{t\geq 0}$ with infinitesimal
generator $L$ acting on a Riemannian manifold $(N, g_N)$, the
$\Gamma_2$-calculus of Bakry--\'Emery~\cite{BakryEmery1985}
defines, for smooth $f, g$ on $N$,
$$
\Gamma(f,g) = \tfrac12\bigl(L(fg) - f L g - g L f\bigr),\qquad
\Gamma_2(f,g) = \tfrac12\bigl(L \Gamma(f,g) - \Gamma(L f, g) - \Gamma(f, L g)\bigr).
$$
The curvature--dimension condition $\mathrm{CD}(K,\infty)$ is the inequality
$\Gamma_2(f,f) \geq K\,\Gamma(f,f)$ for all smooth $f$. For the
Laplace--Beltrami operator on a Riemannian manifold, this is equivalent to
$\Ric_{g_N} \geq K\,g_N$.

\subsection{Babelon--Viallet O'Neill formula}
Babelon--Viallet~\cite[Eq.~(4.5)]{BabelonViallet1981} (see also
Mitter--Viallet~\cite{MitterViallet1981}) gave the O'Neill formula for the
horizontal/vertical decomposition of $T\mathcal A = T\mathcal A^\mathrm{horiz}
\oplus T\mathcal A^\mathrm{vert}$ at $A \in \Omega$, identifying the Ricci
tensor on the quotient $\mathcal A/\mathcal G$ (when restricted to a region
where the vertical sectional curvature is controlled by $\spec(M[A])$):
schematically, $\Ric_{\mathcal A/\mathcal G}$ admits a lower bound of the
form
\begin{equation}\label{eq:BV-ON}
\Ric_{\mathcal A/\mathcal G}(\xi,\xi) \;\geq\;
C_{\mathrm{ON}}\cdot \lambda_{\min}(M[A])\cdot \abs{\xi}^2_{\mathcal A/\mathcal G},
\quad \xi \in T_A(\mathcal A/\mathcal G),
\end{equation}
where $C_{\mathrm{ON}} > 0$ is a constant absorbing factors from the
horizontal projection $\pi$ and the O'Neill bracket (citation to verify;
the precise constant in the original Babelon--Viallet form requires a
careful re-derivation in the language of $H^1_{\Coul}$ on
$\overline\Lambda_{S_0}$ and is fixed by~\eqref{eq:BV-ON} to satisfy the
normalisation $\Ric \geq \lambda_{\min}\cdot g$ when the vertical part
is trivial). We work below with the normalisation $C_{\mathrm{ON}} = 1$,
which is the convention adopted in
\cite[Cor.~KR--FP--A]{RemondiereKRFP3_2026}.

\subsection{LSI and spectral gap}
The classical Rothaus theorem~\cite{Rothaus1981}: if a probability measure
$\mu$ satisfies a logarithmic Sobolev inequality with constant $\rho>0$,
$$
\Ent_\mu(f^2) \leq \tfrac{2}{\rho}\int \abs{\nabla f}^2 d\mu,
$$
then the spectral gap of the associated Dirichlet form satisfies
$\lambda_1 \geq \rho$. Otto--Villani~\cite{OttoVillani2000} gave the converse
implication LSI $\Leftrightarrow$ Talagrand $T_2$ + suitable transport
estimates, and in particular: a $\mathrm{CD}(K,\infty)$ condition with
$K > 0$ implies LSI with constant $K$ for the heat semigroup. Cattiaux--Guillin
and many subsequent authors extended this to the infinite-dimensional Markov
setting.

% =========================================================================
\section{The KR--FP chain and Corollary KR--FP--A}\label{sec:KRFPchain}

We recall here in compact form the three statements of the companion
work~\cite{RemondiereKRFP3_2026}, indexed KR--FP--1, KR--FP--2, KR--FP--3,
plus the immediate Corollary KR--FP--A which is the input to the
Bakry--\'Emery step.

\begin{proposition}[KR--FP--1, Birman--Schwinger / Aubin--Talenti bound]
\label{prop:KRFP1}
Let $A \in H^1_{\Coul}(M; \Omega^1\otimes\mathrm{ad}\,P)$. Write
$M[A] = -\Delta(\mathrm{Id} + \widetilde K(A))$. Then
$$
\norm{\widetilde K(A)}_{\mathcal B(L^2)}
\;\leq\; 2 C_S g\,\norm{A}_{L^4}
\;+\; C_S^2 g^2\,\norm{A}_{L^4}^2,
\qquad C_S = (3/(4\pi^2))^{1/4} \approx 0.392.
$$
\end{proposition}
\begin{proof}
See~\cite[Lemma~4.1 (Birman--Schwinger)]{RemondiereKRFP3_2026}, building on
the optimal Aubin--Talenti Sobolev embedding
$\dot H^1(\R^4) \hookrightarrow L^4(\R^4)$ \cite{Aubin1976,Talenti1976}.
\end{proof}

\begin{proposition}[KR--FP--2, Kostant identity]\label{prop:KRFP2}
For $h$ in a Cartan subalgebra $\hfrak \subset \gfrak$,
$$
\sum_{b=1}^{\dim\gfrak} \norm{[h,T^b]}_\gfrak^2
\;=\; 2\sum_{\alpha\in\Phi^+}\alpha(h)^2
\;=\; \norm{h}_\gfrak^2.
$$
\end{proposition}
\begin{proof}
Kostant~\cite[\S 5]{Kostant1959}; root-space decomposition of
$-\Tr(\ad_h^2)$. See also~\cite[Proposition~2.3 (Kostant)]{RemondiereKRFP3_2026}.
\end{proof}

\begin{theorem}[KR--FP--3, conditional spectral bound on $\overline\Lambda_{S_0}$]
\label{thm:KRFP3}
Conditional on the three hypotheses (H1)~generic-vanishing of minimising
eigenfunctions, (H2)~compact-manifold Sobolev constant uniform within
$1+\varepsilon$, (H3)~measurable Cartan selection on $\overline\Lambda$
(see~\cite[Remark~5.3 (Status of (H1)--(H3))]{RemondiereKRFP3_2026}),
for every
$A \in \overline\Lambda_{S_0}$:
$$
\lambda_{\min}(d_A^\dagger d_A) \;\geq\; m_0^2\,(1-\kFP) \;>\; 0,
$$
where $m_0^2 = \inf\spec(-\Delta\restriction_\mathrm{transverse})$ on
$\gfrak$-valued transverse 1-forms, and
$\kFP = 1/(2|\Phi^+(G)|) = 1/6$ for $G=SU(3)$.
\end{theorem}
\begin{proof}
See~\cite{RemondiereKRFP3_2026}.
\end{proof}

\begin{corollary}[KR--FP--A, Ricci curvature bound on $\overline\Lambda_{S_0}$]
\label{cor:KRFPA}
Under the hypotheses (H1)--(H3) of Theorem~\ref{thm:KRFP3} and with the
normalisation $C_\mathrm{ON} = 1$ in~\eqref{eq:BV-ON},
$$
\Ric_{\mathcal A/\mathcal G}\restriction_{\overline\Lambda_{S_0}}
\;\geq\; (1-\kFP)\,m_0^2 \cdot g_{\mathcal A/\mathcal G}.
$$
\end{corollary}
\begin{proof}
Apply~\eqref{eq:BV-ON} with $\lambda_{\min}(M[A]) \geq (1-\kFP) m_0^2$ from
Theorem~\ref{thm:KRFP3}. Tangential to $\overline\Lambda_{S_0}$ the
metric $g_{\mathcal A/\mathcal G}$ is well defined since
$\overline\Lambda_{S_0}\subset\overline\Omega$, where $M[A]$ is invertible
in the interior.
\end{proof}

\begin{remark}
The numerical value $\kFP = 1/6$ for $SU(3)$ is confirmed by JAX SU(2)
lattice computations on $L\in\{4,8,12,16,20\}$ extrapolating to
$\norm{K}_\mathrm{emp} \to 0.18 \approx 1/6$
(see~\cite[\S 6 (Numerical observations), Caveat~1]{RemondiereKRFP3_2026}
on $SU(2)$ vs $SU(3)$ matching).
\end{remark}

% =========================================================================
\section{Bakry--\'Emery: from $\Ric$ to LSI}\label{sec:bakry}

We now lift Corollary~\ref{cor:KRFPA} from a pointwise Ricci bound on a
Riemannian set to a uniform logarithmic Sobolev inequality on the natural
Markov semigroup associated with the formal Yang--Mills heat flow on
$\overline\Lambda_{S_0}\subset\mathcal A/\mathcal G$.

\subsection{Symbolic Markov semigroup}
On $(\mathcal A/\mathcal G, g_{\mathcal A/\mathcal G})$, the Laplace--Beltrami
operator $\Delta_{\mathcal A/\mathcal G}$ generates a heat semigroup
$P_t = e^{t\Delta_{\mathcal A/\mathcal G}}$. We restrict $P_t$ to the
closed action-bounded slice $\overline\Lambda_{S_0}$ and consider the
killed semigroup $P_t^{\overline\Lambda_{S_0}}$ with reflecting boundary
on $\partial\overline\Lambda_{S_0}\cap\partial\Omega$
(in the regularised sense; see Remark~\ref{rem:boundary} below).

\subsection{Bakry--\'Emery on a manifold with Ricci bound}

\begin{theorem}[Bakry--\'Emery 1985, geometric version]\label{thm:BE}
Let $(N, g_N)$ be a smooth complete Riemannian manifold with Laplace--Beltrami
operator $L = \Delta_{g_N}$. If $\Ric_{g_N} \geq K\, g_N$ for some
$K > 0$, then the heat semigroup $P_t = e^{tL}$ satisfies the curvature--
dimension condition $\mathrm{CD}(K, \infty)$:
$$
\Gamma_2(f,f) \;\geq\; K\,\Gamma(f,f), \qquad \forall f \in C_c^\infty(N).
$$
Consequently the heat-kernel measure $\mu_t = P_t^*\delta_{x_0}$ satisfies a
logarithmic Sobolev inequality with constant $\rho \geq K(1 - e^{-2Kt})/(2K)$,
which in the stationary limit $t\to\infty$ becomes simply $\rho \geq K$.
\end{theorem}
\begin{proof}
Standard $\Gamma$-calculus; see~\cite{BakryEmery1985} and
the textbook~\cite{BakryGentilLedoux2014}. For the LSI version see in
particular~\cite[Thm.~4.7.2]{BakryGentilLedoux2014}.
\end{proof}

\subsection{Application to $\overline\Lambda_{S_0}$}

\begin{proposition}[Symbolic LSI on $\overline\Lambda_{S_0}$]\label{prop:LSI-Lambda}
Conditional on Corollary~\ref{cor:KRFPA}, the heat semigroup
$P_t^{\overline\Lambda_{S_0}}$ on $(\overline\Lambda_{S_0},
g_{\mathcal A/\mathcal G})$ satisfies a logarithmic Sobolev inequality
with constant
$$
\rho \;\geq\; K \,:=\, (1-\kFP)\,m_0^2.
$$
\end{proposition}
\begin{proof}
By Corollary~\ref{cor:KRFPA}, $\Ric \geq K\,g$ on $\overline\Lambda_{S_0}$
with $K = (1-\kFP) m_0^2 > 0$. The Bakry--\'Emery
Theorem~\ref{thm:BE} applies, giving a uniform LSI with constant $\rho \geq K$
for the symbolic semigroup. The reflection boundary condition is compatible
with the LSI in finite-dimensional analogues (Bakry--\'Emery on manifolds
with boundary, see~\cite{Wang2014}); we adopt this as a working hypothesis
for the formal infinite-dimensional case
(see~Remark~\ref{rem:boundary} on regularisation).
\end{proof}

\begin{remark}[Boundary regularisation]\label{rem:boundary}
The set $\overline\Lambda_{S_0}$ has a boundary inside $\mathcal A/\mathcal G$:
the topological boundary is $\partial\overline\Lambda_{S_0}\cap
\partial\Omega \cup \partial\overline\Lambda_{S_0}\cap\{S_\mathrm{YM} = S_0\}$.
On the first component the Faddeev--Popov spectrum closes, on the second
the action achieves its supremum on the slice. Standard Wentzell or reflection
boundary conditions preserve the LSI in finite-dimensional analogues
(see~\cite{Wang2014} for the manifold-with-boundary version of
Bakry--\'Emery, and~\cite[Sec.~5]{BakryGentilLedoux2014}). We adopt this as a
working hypothesis. A rigorous treatment of the infinite-dimensional boundary
regularisation is part of the input (ii) below (the BBD program).
\end{remark}

\subsection{Otto--Villani transport equivalence}

\begin{proposition}[Otto--Villani 2000, transport-LSI equivalence]
\label{prop:OV}
Conditional on $\mathrm{CD}(K,\infty)$ with $K > 0$, the measure $\mu_\infty
= \lim_{t\to\infty} P_t^*\delta_{x_0}$ (when it exists) satisfies
both Talagrand's $T_2$ transport inequality with constant $K$ and the
logarithmic Sobolev inequality with constant $K$.
\end{proposition}
\begin{proof}
Otto--Villani~\cite{OttoVillani2000}; see also the streamlined treatment
in~\cite[Sec.~3]{BakryGentilLedoux2014}.
\end{proof}

\paragraph{Net output of Section~\ref{sec:bakry}.}
Under hypotheses (H1)--(H3) of Theorem~\ref{thm:KRFP3}, the symbolic Markov
semigroup on $\overline\Lambda_{S_0}\subset\mathcal A/\mathcal G$ satisfies
a uniform LSI with constant $\rho \geq (1-\kFP)m_0^2 > 0$.

% =========================================================================
\section{Continuum descent}\label{sec:descent}

The output of Section~\ref{sec:bakry} is a uniform LSI for the formal
\emph{geometric} heat flow on a closed slice of the orbit space, with
constant independent of any lattice approximation. To convert this to a
mass gap for the physical Yang--Mills measure
$\mu_\mathrm{YM} = Z^{-1}\exp(-S_\mathrm{YM}/g^2)\,\mathcal D A$ (with the
appropriate regularisation), we must descend through the lattice
approximations and demonstrate that the LSI constant survives the limit.

\subsection{Lattice approximation}
Let $\Lambda_a \subset (a\Z)^4$ be a finite lattice of spacing $a > 0$ and
side $L_a$. The Wilson measure
$\mu_a = Z_a^{-1}\exp(-S_{\mathrm{W},a}[U]/g^2)\,d U$ with
$S_{\mathrm{W},a}[U] = \sum_p (1 - \tfrac{1}{N}\Re\Tr U_p)$ defines a
$G^{\#\mathrm{links}}$-valued probability measure. Heat-bath / Langevin
dynamics on $\mu_a$ are well-studied (Parisi--Wu stochastic
quantisation; Brownian motion on $G^{\#\mathrm{links}}$).

\subsection{The BBD uniform LSI input}\label{sec:BBD}

The principal external input of this Letter, beyond
Theorem~\ref{thm:KRFP3} and the Bakry--\'Emery step, is the following
external assumption, modelled on the scalar
$\varphi^4$-result~\cite{BauerschmidtDagallier2024}.

\begin{hypothesis}[Uniform LSI for 4D pure $SU(N)$ Wilson measure]
\label{hyp:BBD}
There exists a constant $\cinf(D) > 0$, depending only on dimension $D = 4$
and on $G = SU(N)$, such that for every $a > 0$ and every $L_a$ in the
appropriate scaling window, the Wilson measure $\mu_a$ on the periodic
$\Lambda_a$ satisfies a logarithmic Sobolev inequality with constant
$\rho_a \geq 1/\cinf(D)$. Equivalently, the spectral gap of the
heat-bath / Langevin generator on $\mu_a$ is at least $1/\cinf(D)$,
uniformly in $a$.
\end{hypothesis}

\paragraph{Status of Hypothesis~\ref{hyp:BBD}.} For scalar
$\varphi^4$-measures in dimensions $d = 2, 3$ this is exactly the result
of Bauerschmidt--Dagallier~\cite{BauerschmidtDagallier2024} via the Polchinski
equation method of \cite{BauerschmidtBodineauDagallier2024}. For 4D
$SU(N)$ pure gauge it is \emph{open}; its derivation is the standard
cluster-expansion problem in the modern probabilistic formulation. The
Cao--Park--Sheffield~\cite{CaoParkSheffield2024} and
Cao--Nissim--Sheffield~\cite{CaoNissimSheffield2025} works on lattice
Yang--Mills as random surfaces, together with the Chandra--Chevyrev--Hairer--Shen
3D construction~\cite{ChandraChevyrevHairerShen2024}, provide important
recent progress but do not yet yield Hypothesis~\ref{hyp:BBD}
in 4D.

\subsection{Dirichlet form descent}
Granted Hypothesis~\ref{hyp:BBD}, the standard Fukushima--Oshima--Takeda
machinery~\cite{FukushimaOshimaTakeda1994} for Dirichlet forms ensures that
the family $(\mu_a)$ is tight in the appropriate distributional space, and
that the limit point $\mu_\infty$ inherits the LSI with constant
$1/\cinf(D)$ (by Mosco-type lower semicontinuity of LSI under regular
graph convergence; see~\cite[Ch.~4]{BakryGentilLedoux2014}).

\subsection{Coupling KR--FP--A to Hypothesis~\ref{hyp:BBD}}
Two LSIs are now in play: the symbolic Bakry--\'Emery LSI from
Proposition~\ref{prop:LSI-Lambda} with constant $(1-\kFP)m_0^2$ on
$\overline\Lambda_{S_0}$, and the BBD LSI with constant $1/\cinf(D)$ on
the physical Wilson measure. The coupling rests on the heuristic identification,
correct in the formal continuum limit, between the symbolic heat flow on
$\mathcal A/\mathcal G\restriction_{\overline\Lambda_{S_0}}$ and the
gauge-invariant component of the Wilson measure restricted to the analogous
Coulomb-gauge slice. We \emph{conjecture} (and the Bauerschmidt collaboration
roadmap of Section~\ref{sec:honest} aims to formalise) that the effective
LSI constant on the gauge-fixed Wilson measure is
$$
\rho_\mathrm{eff} \;\geq\; \frac{(1-\kFP)\,m_0^2}{\cinf(D)},
$$
i.e.\ the product of the symbolic LSI with the BBD LSI \emph{divided by
$\cinf(D)$} (the standard cluster expansion contribution). This is the
input that produces~\eqref{eq:main-informal}.

\begin{remark}[Multiplicative structure of the constant]
The factor $1/\cinf(D)$ arises because the BBD LSI gives the lattice
$\mu_a$ a normalised spectral gap $1/\cinf(D)$, and the geometric
component $(1-\kFP) m_0^2$ from Proposition~\ref{prop:LSI-Lambda}
\emph{multiplies} this in the saturated Bakry--\'Emery scheme of
\cite[Sec.~IV.B]{RemondiereMassGapPRL_2026}, where the explicit
cross-group formula
$\mathrm{C}_\mathrm{LSI}(G,D) = \cinf(D)\cdot f(\pi_1(G))\cdot
[1 - \kFP\,\delta_\mathrm{sat}]$ has been derived empirically and
shown to match SU(N) lattice gauge data with mean residual $5$--$8\%$
across $N = 2, 3, 4, 5$ in $D = 3, 4, 5, 6$
\cite[Tab.~III]{RemondiereMassGapPRL_2026}.
The empirical $\cinf(4) = 0.247$ from Wilson-flow plateau measurements
(\cite[Sec.~V]{RemondiereMassGapPRL_2026}) gives, in the SU(3) case,
$\Delta_\mathrm{est} \geq (5/6)\cdot m_0^2 / 0.247 \approx 3.38\,m_0^2$.
\end{remark}

% =========================================================================
\section{Main theorem (conditional)}\label{sec:thm}

We now state the conditional reduction in its rigorous form.

\begin{theorem}[KR--FP--B, mass gap conditional reduction]\label{thm:main}
Let $G$ be a compact simple Lie group with Lie algebra $\gfrak$ and
positive roots $\Phi^+$. Set $\kFP := 1/(2|\Phi^+|)$.
Assume:
\begin{enumerate}[label={\rm(\roman*)},leftmargin=2.5em]
\item Hypotheses {\rm (H1), (H2), (H3)} of Theorem~\ref{thm:KRFP3}
{\rm \cite{RemondiereKRFP3_2026}}, i.e.\ the Faddeev--Popov spectral bound
$\lambda_{\min}(d_A^\dagger d_A) \geq m_0^2(1-\kFP)$ on
$\overline\Lambda_{S_0}\subset\overline\Lambda \subset \mathcal A/\mathcal G$;
\item Hypothesis~\ref{hyp:BBD}: a uniform logarithmic Sobolev inequality with
constant $1/\cinf(D)$ for the 4D pure $SU(N)$ Wilson measure on the
periodic lattice $\Lambda_a$, uniformly in $a > 0$ across the scaling window;
\item Compatibility {\rm (C)} of the symbolic and BBD LSIs in the sense that
the gauge-fixed projection of the Wilson measure on
$\overline\Lambda_{S_0}$ inherits both LSIs, and the resulting effective
LSI constant on the orbit-space measure is at least
$(1-\kFP)\,m_0^2/\cinf(D)$ in the continuum limit.
\end{enumerate}
Then the spectral gap $\Delta$ of the 4D pure $SU(N)$ Yang--Mills
Hamiltonian satisfies
\begin{equation}\label{eq:main}
\Delta \;\geq\; \frac{(1-\kFP)\,m_0^2}{\cinf(D)} \;>\; 0.
\end{equation}
In particular for $G = SU(3)$, $D = 4$,
$$
\Delta \;\geq\; \tfrac{5}{6}\cdot\frac{m_0^2}{\cinf(4)} \;>\; 0.
$$
\end{theorem}

\begin{proof}
Step 1: By (i) and Theorem~\ref{thm:KRFP3}, the Faddeev--Popov spectrum
on $\overline\Lambda_{S_0}$ is bounded below by $(1-\kFP)m_0^2$.

Step 2: By Corollary~\ref{cor:KRFPA} and the Babelon--Viallet O'Neill
formula~\eqref{eq:BV-ON}, $\Ric_{\mathcal A/\mathcal G} \geq (1-\kFP)
m_0^2 \cdot g$ on $\overline\Lambda_{S_0}$.

Step 3: By Theorem~\ref{thm:BE} (Bakry--\'Emery), the symbolic Markov
semigroup on $(\overline\Lambda_{S_0},g_{\mathcal A/\mathcal G})$ satisfies
$\mathrm{CD}((1-\kFP)m_0^2,\infty)$.

Step 4: By Proposition~\ref{prop:OV} (Otto--Villani) and
Proposition~\ref{prop:LSI-Lambda}, the heat-kernel measure on
$\overline\Lambda_{S_0}$ obeys a logarithmic Sobolev inequality with
constant $\rho \geq (1-\kFP)m_0^2$.

Step 5: By Rothaus~\cite{Rothaus1981}, the symbolic LSI implies a
spectral gap $\lambda_1 \geq (1-\kFP)m_0^2$ for the symbolic Dirichlet form
on $\overline\Lambda_{S_0}$.

Step 6: By (ii), Hypothesis~\ref{hyp:BBD}, the lattice Wilson measure
$\mu_a$ has uniform LSI with constant $1/\cinf(D)$, and the
Fukushima--Oshima--Takeda Dirichlet-form Mosco-convergence
machinery~\cite{FukushimaOshimaTakeda1994} delivers the continuum measure
$\mu_\infty$ with the same LSI constant.

Step 7: By (iii), the compatibility hypothesis, the effective LSI constant
on the orbit-space continuum measure $\mu_\infty^{\mathcal A/\mathcal G}$
is at least $(1-\kFP)m_0^2/\cinf(D)$.

Step 8: By the Reed--Simon transfer between LSI and the spectral gap of the
associated Markov semigroup (an immediate consequence of
Rothaus~\cite{Rothaus1981} applied to $\mu_\infty^{\mathcal A/\mathcal G}$),
the spectral gap of the physical Yang--Mills Hamiltonian satisfies
$\Delta \geq (1-\kFP)m_0^2/\cinf(D) > 0$.

For $G = SU(3)$: $\kFP = 1/6$, $1-\kFP = 5/6$.
\end{proof}

\begin{remark}[The compatibility hypothesis (C) is not automatic]
Compatibility (C) is the meeting point of the geometric and probabilistic
analyses. It encodes the assertion that the projection of the lattice
Wilson measure onto a Coulomb-gauge slice of the orbit space inherits
\emph{both} the geometric LSI (with the constant from
Proposition~\ref{prop:LSI-Lambda}) and the BBD LSI (with the constant from
Hypothesis~\ref{hyp:BBD}). Empirically (cf.\
\cite[Tab.~III]{RemondiereMassGapPRL_2026}) the resulting LSI constant
\emph{multiplies} the two factors, suggesting that the Cartan-saturation
correction from KR--FP--3 acts as a structural multiplicative shift on
the BBD baseline $\cinf(D)$. A rigorous proof of (C) would require a
gauge-fixing argument in the BBD formalism, which we leave to future work.
\end{remark}

\begin{remark}[Numerical illustration for $G=SU(3)$, $D=4$]
With the empirical Wilson-flow plateau value
$\cinf^{\mathrm{emp}}(4) = 0.247\pm 0.001$
(\cite[Tab.~II]{RemondiereMassGapPRL_2026}), the conditional bound reads
$$
\Delta \;\geq\; \tfrac{5}{6}\cdot \frac{m_0^2}{0.247}
\;\approx\; 3.38\, m_0^2.
$$
With $m_0^2 \sim \Lambda_{\mathrm{QCD}}^2 \sim (250\,\mathrm{MeV})^2$
this gives $\Delta \gtrsim (460\,\mathrm{MeV})^2$, of the correct order of
magnitude as the lattice glueball mass $m(0^{++})_{SU(3)} \simeq 1700$~MeV
(Morningstar--Peardon~\cite{Morningstar1999}). The agreement should not be
overstated: the absolute prefactor $m_0^2$ is unknown
\textit{a priori}, and only the dependence on $\kFP$ and $\cinf(D)$ is
delivered by Theorem~\ref{thm:main}.
\end{remark}

% =========================================================================
\section{Honest conditionality and formalisation roadmap}\label{sec:honest}

We list explicitly the remaining open problems and provide a tentative
roadmap.

\subsection{Open hypotheses}

\begin{itemize}[leftmargin=2em]
\item \textbf{(H1)~Generic-vanishing} \cite{RemondiereKRFP3_2026}.
The asymptotic alignment of minimising eigenfunctions
$\eta_\mathrm{min}(A)$ of $M[A]$ with the Cartan subspace
$\hfrak\otimes\Omega^1_\mathrm{transverse}$ near $\partial\Omega$.
Numerical evidence in $SU(2)$ lattice computations
\cite[Table~1 ($SU(2)$ lattice)]{RemondiereKRFP3_2026} (Caveat~1: matching
to $\kFP = 1/6$ rather than $\kFP^{SU(2)} = 1/2$ requires further analysis).
Theoretical proof requires perturbation analysis on $\overline\Lambda_{S_0}$
near the horizon; estimated effort 4--7 months full-time.

\item \textbf{(H2)~Compact-manifold Sobolev constant}. The optimal
Aubin--Talenti constant $C_S = (3/(4\pi^2))^{1/4}$ is for the unbounded
$\R^4$. On the compact $M$ a quantitative lower-order correction is needed,
depending on $\mathrm{vol}(M)$, $\mathrm{inj}(M)$, $\mathrm{Ric}(M)$.
This is a standard but unwritten estimate. Effort: 1--2 months.

\item \textbf{(H3)~Measurable Cartan selection on $\overline\Lambda$}.
Standard at regular points; care required at degeneracy strata. Hahn--Banach
plus measurable selection theorems should suffice. Effort: 1 month.

\item \textbf{Hypothesis~\ref{hyp:BBD}~(BBD uniform LSI for 4D pure $SU(N)$)}.
The principal external input. The scalar version is in
\cite{BauerschmidtDagallier2024}; the non-Abelian gauge adaptation
involves substantial technical work (Polchinski equation for $SU(N)$ Haar
measure, cluster expansion control). Estimated effort with Bauerschmidt
collaboration: 18--24 months. A first email contact has been initiated
by the author (May 2026).

\item \textbf{Compatibility (C)}. Coupling of geometric and BBD LSIs via
a gauge-fixed projection. Conditional on the gauge-fixing being implemented
in the Polchinski-RG framework. Effort: 2--4 months conditional on the
BBD-collaboration framework.
\end{itemize}

\subsection{Path probability estimate}
Following the convention of~\cite[Tab.~IV]{RemondiereMassGapPRL_2026}, we
quantify the probability that the full chain closes within 10 years as
follows: under the Bauerschmidt-collaboration path, conditional on
(H1)--(H3) and Hypothesis~\ref{hyp:BBD} all being discharged,
$P(\text{Theorem~\ref{thm:main} formalised within 10y}) \approx 45$--$60\%$,
i.e.\ the dominant uncertainty is on Hypothesis~\ref{hyp:BBD} for $SU(N)$
and on the compatibility (C). Conditional on the direct path (H1)--(H3) only
(no collaboration), $P \approx 35$--$50\%$ at the price of 4--7 months on
(H1) alone.

\subsection{Strict bypass note}
The reduction articulated here does \emph{not} bypass the
Ba\l aban--Bauerschmidt cluster-expansion programme. It bypasses the
cluster expansion \emph{for the Faddeev--Popov spectral lower bound}
(Theorem~\ref{thm:KRFP3} replaces it with a geometric/Lie-algebraic
argument), but still requires the cluster expansion (in the BBD form,
Hypothesis~\ref{hyp:BBD}) for the measure descent. This is a partial bypass.

% =========================================================================
\section{Conclusion}\label{sec:conclusion}

We have articulated a conditional reduction of the 4D pure $SU(N)$ Yang--Mills
mass gap problem along the chain
$$
\text{KR--FP--3}
\;\Rightarrow\;\text{KR--FP--A (Babelon--Viallet O'Neill)}
\;\Rightarrow\;\text{Bakry--\'Emery}
\;\Rightarrow\;\text{Otto--Villani LSI}
\;\Rightarrow\;\text{Rothaus gap}
\;\Rightarrow\;\text{continuum gap (BBD + Fukushima--Oshima--Takeda)}.
$$
The output is the bound
$\Delta \geq (1-\kFP)\,m_0^2/\cinf(D)$, which for $G = SU(3)$ in $D = 4$
reads $\Delta \geq \tfrac{5}{6}\,m_0^2/\cinf(4)$.
The contribution is structural: a single conceptual chain that exhibits
explicitly which open analytic problems remain. These are the three KR--FP--3
hypotheses (H1)--(H3) of~\cite{RemondiereKRFP3_2026}, the BBD-style uniform
LSI Hypothesis~\ref{hyp:BBD} for 4D pure $SU(N)$ gauge in the spirit
of~\cite{BauerschmidtBodineauDagallier2024,BauerschmidtDagallier2024}, and
the compatibility (C). We do \emph{not} claim a proof of the Clay mass gap.

\subsection*{Acknowledgments}
The author thanks anonymous referees of the companion KR--FP--3 manuscript
for criticism that sharpened the formulation of the three hypotheses.
The lattice computations supporting the numerical verification of the
geometric route were performed on an NVIDIA RTX 5060 Ti GPU. This work is
released under a CC-BY-4.0 license.

% =========================================================================
\begin{thebibliography}{99}

\bibitem{Aubin1976}
T.~Aubin,
\emph{Probl\`emes isop\'erim\'etriques et espaces de Sobolev},
J.~Differential Geom.~\textbf{11} (1976), 573--598.

\bibitem{BabelonViallet1981}
O.~Babelon and C.~M.~Viallet,
\emph{The Riemannian geometry of the configuration space of gauge theories},
Comm.~Math.~Phys.~\textbf{81} (1981), 515--525.

\bibitem{BakryEmery1985}
D.~Bakry and M.~\'Emery,
\emph{Diffusions hypercontractives},
in \emph{S\'eminaire de Probabilit\'es XIX 1983/84}, Lecture Notes in Math.
\textbf{1123}, Springer (1985), 177--206.

\bibitem{BakryGentilLedoux2014}
D.~Bakry, I.~Gentil, and M.~Ledoux,
\emph{Analysis and Geometry of Markov Diffusion Operators},
Grundlehren der Mathematischen Wissenschaften \textbf{348},
Springer (2014).

\bibitem{Balaban1985}
T.~Ba\l aban,
\emph{Renormalization group approach to lattice gauge field theories.
I.~Generation of effective actions in a small field approximation},
Comm.~Math.~Phys.~\textbf{109} (1987), 249--301.

\bibitem{BauerschmidtBodineauDagallier2024}
R.~Bauerschmidt, T.~Bodineau, and B.~Dagallier,
\emph{Stochastic dynamics and the Polchinski equation: an introduction},
Probab.~Surveys \textbf{21} (2024), 200--290;
\texttt{arXiv:2307.07619}.

\bibitem{BauerschmidtDagallier2024}
R.~Bauerschmidt and B.~Dagallier,
\emph{Log-Sobolev inequality for the $\varphi^4_2$ and $\varphi^4_3$ measures},
Comm.~Pure Appl.~Math.~\textbf{77} (2024), 2579--2612;
\texttt{arXiv:2202.02295}.

\bibitem{CaoNissimSheffield2025}
S.~Cao, R.~Nissim, and S.~Sheffield,
\emph{Dynamical approach to area law for lattice Yang--Mills} (2025);
\texttt{arXiv:2509.04688}.

\bibitem{CaoParkSheffield2024}
S.~Cao, J.~Park, and S.~Sheffield,
\emph{Random surfaces and lattice Yang--Mills},
Commun.~Amer.~Math.~Soc.~(accepted, 2024);
\texttt{arXiv:2307.06790}.

\bibitem{ChandraChevyrevHairerShen2024}
A.~Chandra, I.~Chevyrev, M.~Hairer, and H.~Shen,
\emph{Stochastic quantisation of Yang--Mills--Higgs in 3D},
Invent.~Math.~(2024);
\texttt{arXiv:2201.03487}.

\bibitem{Chatterjee2024}
S.~Chatterjee,
\emph{A scaling limit of $SU(2)$ lattice Yang--Mills--Higgs theory},
Probab.~Math.~Phys.~(accepted, 2024);
\texttt{arXiv:2401.10507}.

\bibitem{DellAntonioZwanziger1991}
G.~Dell'Antonio and D.~Zwanziger,
\emph{Every gauge orbit passes inside the Gribov horizon},
Comm.~Math.~Phys.~\textbf{138} (1991), 291--299.

\bibitem{FukushimaOshimaTakeda1994}
M.~Fukushima, Y.~\=Oshima, and M.~Takeda,
\emph{Dirichlet Forms and Symmetric Markov Processes},
de Gruyter Studies in Mathematics \textbf{19}, Walter de Gruyter (1994).

\bibitem{Gribov1978}
V.~N.~Gribov,
\emph{Quantization of nonabelian gauge theories},
Nuclear Phys.~B \textbf{139} (1978), 1--19.

\bibitem{JaffeWitten2000}
A.~Jaffe and E.~Witten,
\emph{Quantum Yang--Mills theory}, Clay Mathematics Institute Millennium
Problem description (2000).

\bibitem{Kostant1959}
B.~Kostant,
\emph{The principal three-dimensional subgroup and the Betti numbers of a
complex simple Lie group},
Amer.~J.~Math.~\textbf{81} (1959), 973--1032.

\bibitem{MitterViallet1981}
P.~K.~Mitter and C.~M.~Viallet,
\emph{On the bundle of connections and the gauge orbit manifold in
Yang--Mills theory},
Comm.~Math.~Phys.~\textbf{79} (1981), 457--472.

\bibitem{Morningstar1999}
C.~J.~Morningstar and M.~J.~Peardon,
\emph{The glueball spectrum from an anisotropic lattice study},
Phys.~Rev.~D \textbf{60} (1999), 034509;
\texttt{arXiv:hep-lat/9901004}.

\bibitem{OttoVillani2000}
F.~Otto and C.~Villani,
\emph{Generalization of an inequality by Talagrand and links with the
logarithmic Sobolev inequality},
J.~Funct.~Anal.~\textbf{173} (2000), 361--400.

\bibitem{RemondiereKRFP3_2026}
K.~R\'emondi\`ere,
\emph{Conditional spectral bound for the Faddeev--Popov operator on the
fundamental modular domain via Lie-algebraic reduction (KR--FP--3)},
preprint, Annals of Mathematics submission, May 2026.

\bibitem{RemondiereMassGapPRL_2026}
K.~R\'emondi\`ere,
\emph{Mass gap from first principles: a cross-group formula and conditional
reduction},
preprint, Physical Review Letters submission, May 2026.

\bibitem{Rothaus1981}
O.~S.~Rothaus,
\emph{Diffusion on compact Riemannian manifolds and logarithmic Sobolev
inequalities},
J.~Funct.~Anal.~\textbf{42} (1981), 102--109.

\bibitem{Singer1978}
I.~M.~Singer,
\emph{Some remarks on the {Gribov} ambiguity},
Comm.~Math.~Phys.~\textbf{60} (1978), 7--12.

\bibitem{Talenti1976}
G.~Talenti,
\emph{Best constant in {Sobolev} inequality},
Ann.~Mat.~Pura Appl.~\textbf{110} (1976), 353--372.

\bibitem{Uhlenbeck1982}
K.~K.~Uhlenbeck,
\emph{Connections with $L^p$ bounds on curvature},
Comm.~Math.~Phys.~\textbf{83} (1982), 31--42.

\bibitem{Wang2014}
F.-Y.~Wang,
\emph{Analysis for Diffusion Processes on Riemannian Manifolds},
Advanced Series on Statistical Science \& Applied Probability \textbf{18},
World Scientific (2014).

\end{thebibliography}

\end{document}
